\section{Threats to Validity}
\label{sec:threats}

\textbf{One corpus, one domain.} Every measurement comes from a single snapshot of
one filing system. SEC disclosure is unusually favourable to structural retrieval:
it is mandated, schematized, and carries a stable public entity identifier in the
CIK. Corpora without an authoritative identifier --- internal wikis, support
tickets, email --- would require the entity-resolution step this paper largely
assumes, and \Cref{sec:failures} shows that step is our weakest layer. Nothing
here should be read as generalizing to such corpora.

\textbf{Twenty questions is a small sample.} Each category rests on five
questions, so a category-level rate has a wide confidence interval: a single
question changes a category by 20 percentage points at question level and 3
attempts out of 15 at attempt level. The category results are directional, not
conclusive. The one result we would defend as strong is T3, where the margin is
15 of 15 against 0 of 15 for both baselines and the mechanism is independently
diagnosed.

\textbf{The questions were written by the author of the graph.} This is the most
serious threat and it cannot be fully mitigated. We took three measures: two of
the four categories were pre-registered as ones the graph should lose, and it did
lose both; the T3 subjects were selected by a measured criterion rather than by
inspection of results (\Cref{sec:design}); and the ground truth was verified
against the loaded tables rather than against model outputs. None of this
eliminates the risk that the question set favours the design in ways we cannot
see. An independent question set is the obvious next step.

\textbf{One model, and the verdict moves with it.} All reported numbers come from
\texttt{deepseek-v4-flash}. The probe in \Cref{tab:crossmodel} shows the featured
question passing on one model, failing on another and being refused by a third
with the retrieval architecture held constant. Our results characterize
\emph{architecture $\times$ model}, and we did not vary the model across the full
180 attempts.

\textbf{The graph inherits EX-21 incompleteness.} Rule 601(b)(21)(ii) allows
filers to omit insignificant subsidiaries, so \texttt{:SUBSIDIARY\_OF} is a floor
on corporate structure. Any question whose correct answer depends on a subsidiary
the filer chose not to disclose is unanswerable from this corpus by any of the
three arms, and our ground truth is defined against what was disclosed rather than
against reality.

\textbf{Category results move with implementation details.} While calibrating
the harness we found an undocumented name-normalisation convention in the schema
descriptions; documenting it changed three of the four categories
(\Cref{sec:harness}). Nothing about any retrieval architecture changed. If one
sentence of schema prose was worth three categories, other such sensitivities are
plausibly latent in results we have not thought to re-examine, and the ones we
would be least likely to notice are those that flatter the design. The same logic
applies to the wins: some part of GraphRAG's margins may reflect implementation
effort we spent on the graph and not on the baselines.
\texttt{:RESOLVES\_TO} coverage of 0.12\% remains a hard bound on what any
entity-resolution question can measure here.

\textbf{Repair asymmetry favours the graph.} As recorded in \Cref{sec:design},
the zero-row repair trigger fires for a graph query that misses but not for an
over-broad SQL join that returns many wrong rows. GraphRAG used its repair on 34
of 60 attempts against text-to-SQL's 7, so it received materially more effective
assistance from an identical rule.

\textbf{A single year, and a moving corpus.} The window spans twelve monthly
feeds. Annual filings cluster seasonally (\Cref{fig:months}), so the mix of form
types is not stationary within it, and the absolute counts are not reproducible
against a later snapshot because EDGAR accumulates.

\textbf{Undetected errors are likely.} We found several defects in our own
harness, one of them a ground-truth error discovered after we believed the
results were final (\Cref{sec:harness}). That we found it late is direct evidence
that others may remain, and we have no way to bound how many.
